% Bagian: computability
% Bab: computability-theory
% Subbagian: def-functions-self-reference

\documentclass[../../../include/open-logic-section]{subfiles}

\begin{document}

\olfileid[id]{cmp}{thy}{slf}
\olsection{Mendefinisikan Fungsi dengan Acuan Diri}

Secara umum, kemampuan mendefinisikan fungsi berdasarkan fungsi itu sendiri
berguna. Sebagai contoh, diberikan fungsi dapat dihitung $k$, $l$, dan~$m$,
lemma titik tetap mengatakan bahwa terdapat fungsi dapat dihitung
parsial~$f$ yang memenuhi persamaan berikut bagi setiap~$y$:
\[
f(y) \simeq
\begin{cases}
k(y) & \text{jika $l(y) = 0$} \\
f(m(y)) & \text{selain itu.}
\end{cases}
\]
Secara lebih khusus, sekali lagi, $f$ diperoleh dengan menetapkan
\[
g(x,y) \simeq
\begin{cases}
k(y) & \text{jika $l(y) = 0$} \\
\cfind{x}(m(y)) & \text{selain itu}
\end{cases}
\]
lalu menggunakan lemma titik tetap untuk menemukan indeks~$e$ sedemikian
sehingga $\cfind{e}(y)\simeq g(e,y)$.

Sebagai contoh konkret, fungsi ``faktor persekutuan terbesar''
$\fn{gcd}(u,v)$ dapat didefinisikan dengan
\[
\fn{gcd}(u,v) \simeq
\begin{cases}
v & \text{jika $u = 0$} \\
\fn{gcd}(\fn{mod}(v, u), u) & \text{selain itu}
\end{cases}
\]
dengan $\fn{mod}(v,u)$ menyatakan sisa pembagian $v$ oleh~$u$. Penerapan
lemma titik tetap menunjukkan bahwa $\fn{gcd}$ dapat dihitung parsial.
(Sebenarnya, ini dapat disajikan dalam bentuk di atas dengan membiarkan $y$
mengodekan pasangan $\tuple{u,v}$.) Induksi berikutnya pada~$u$ lalu
menunjukkan bahwa $\fn{gcd}$ sebenarnya total.

Tentu saja, kita dapat menyusun definisi beracuan-diri yang jauh lebih rumit
daripada contoh-contoh yang baru saja dibahas. Kebanyakan bahasa pemrograman
mendukung definisi fungsi berdasarkan fungsi itu sendiri dengan satu atau
lain cara. Perhatikan bahwa hal ini sedikit kurang dramatis daripada kemampuan
mendefinisikan fungsi berdasarkan suatu \emph{indeks} bagi algoritma yang
menghitung fungsi tersebut, dan kemampuan terakhir inilah yang diberikan
teorema titik tetap dalam keumuman penuhnya.

\end{document}
